\chapter[Ordinary People and Totalitarianism]{How Totalitarian Politics Draws Ordinary People In}
\section{A Radio in a Wooden Case}
First, pick up an object: an old radio in a wooden case.

About half a meter tall, it has a walnut shell and a cloth-covered front concealing its speaker. Two knobs sit at the lower right, one for power and volume, the other for tuning. Switch it on and wait half a minute. Vacuum tubes warm slowly, glowing orange, before sound seeps through the cloth with a faint underlying hiss.

In 1930s Germany, such radios were extraordinarily cheap, affordable even to ordinary workers. The market had not achieved this unaided: government subsidies deliberately made it possible. Its model name was plain: "People's Receiver." Millions sold within a few years, making it one of the fastest-spreading household appliances. Before war began, most German living rooms contained such a glowing wooden box. People nicknamed it, roughly, "Goebbels's loudmouth," after Nazi propaganda minister Joseph Goebbels.

Consider how unusual this was. Newspapers cost money, required literacy, and had to be bought outside; radio did not. Sit at home, turn a knob, and a voice enters your living room: news, speeches, marches, weather, all behind the same cloth. Whoever controlled that sound channel no longer needed to knock at every door daily.

This warm wooden box is our entrance. The chapter asks not "what was broadcast?" but an uncomfortable question: when one voice sounds in your living room every day, and your family merely listens, nods, follows instructions, and stays silent, are you participating in what happens next?

\section{A New Sign in the Bakery Window}
Come into a scene with me.

Time: an early autumn morning in 1935. Place: a small central German town, a cobbled shopping street, a corner bakery. Oven light glows deep inside; freshly baked rye bread gives the air a sour, toasted fragrance. A brass bell rings whenever someone enters.

The baker is about fifty, his apron covered in flour. He has baked here for thirty years and recognizes nearly every face on the street. Before opening today, he does something unprecedented: puts a new paper sign inside the window. The trade association has distributed it to all members, and many shops already display it. One line reads: Jews are not served here.

An elderly woman in a headscarf stands in the queue with a shopping net. When her turn comes, she sees the sign. She has bought black bread here for thirty years and arranged relatives' orders for celebration cakes when the baker married. She says nothing, stands a few seconds, then leaves. The bell rings once.

Everyone behind her sees it. A young man says loudly, "About time." A housewife carrying a basket lowers her head and watches her shoes. When her turn comes, she buys bread in a quieter voice than usual. The baker hands it over without meeting her eyes. Across the street, the shoemaker smokes in his doorway, sees everything, says nothing, stubs out the cigarette, and goes inside.

Remember this street. Nobody strikes anyone; no blood flows; no uniform appears. Yet something has changed: a standardized sign, a raised voice, lowered heads, a departing figure. Three years later, on November 9, 1938, windows along the same row will be smashed as Jewish shops, synagogues, and homes are looted and burned during Kristallnacht. By then, most neighbors will have practiced for three years how to lower their heads, walk around, and pretend not to see.

\section{The Real Question: Demons or Cogs?}
First establish the factual framework, then let the question emerge.

Germany lost the First World War in 1918; its empire collapsed and the Weimar Republic arose amid chaos. The 1919 peace treaty imposed war guilt and enormous reparations. In 1923, hyperinflation made marks wastepaper: a dollar bought trillions, wheelbarrows of banknotes bought bread, and lifetime savings vanished within days. The Great Depression struck in 1929. By 1932, registered unemployment approached six million, roughly one worker in three, excluding those uncounted.

Hitler was appointed chancellor in January 1933. In late March, parliament passed the Enabling Act, transferring legislative power to his cabinet: democracy ended itself procedurally. The Nuremberg Laws of 1935 downgraded Jews from citizens to "subjects." War began in 1939. In January 1942, fifteen senior officials met at a villa in Wannsee, Berlin, to coordinate the "Final Solution" for Europe's Jews. By 1945, approximately six million Jews had been murdered, alongside Roma, disabled people, Polish and Soviet prisoners of war and civilians, political prisoners, and religious dissenters.

Here is the conflict. A few hundred demons did not operate this killing machine. It needed population-register clerks, railway dispatchers, accountants confiscating property, police guarding gates, teachers teaching "racial science," officials sending trade-association notices, the bakery's sign, and neighbors lowering their heads in line. Millions contributed directly or indirectly. Most were ordinary people like you and me: working, supporting families, avoiding trouble, following the crowd.

Our question is no longer the easy "why were they so wicked?" It becomes three increasingly difficult questions. First, if most were neither mad nor fanatical, how did they enter step by step? Second, was this a specifically "German" affliction or a human one? Third, hardest of all, what makes us certain that, born in that time on that street, we would not have lowered our heads in the queue?

\section{Different Voices: Berlin to Moscow}
People occupied different positions in this machine. Invite several voices to speak fully.

\textbf{Victims.} Correct one familiar picture first: victims were not a line of silent shadows. Amid hunger and fear in Warsaw's Jewish ghetto, people ran underground schools and newspapers. Historian Emanuel Ringelblum organized a secret archive, collecting testimonies, leaflets, recipes, and children's essays into metal containers buried underground. They knew they might not survive, yet wanted the future to know what happened. Later generations recovered the archive from the ruins. Even in the machine's depths, people recorded, judged, and resisted rather than merely waited.

\textbf{Perpetrators.} Jerusalem, 1961: Nazi war criminal Adolf Eichmann stands trial inside bulletproof glass. He had organized rail transport of Jews from across Europe to extermination camps. His defense was broadly: I was merely a civil servant obeying orders; decisions were not mine; I handled transport schedules. I hated nobody; I did my duty. Notice its structure: cut an entire chain of killing into small segments, then point to yours and say, "Look, this part has no blood."

\textbf{Bystanders.} After the war, countless ordinary Germans said, "We did not know." Not entirely false: extermination camps stood in occupied Poland, and details were closely guarded. Not entirely true either: noisy arrests of neighbors, colleagues who never returned, closed freight trains, circulating rumors. Everyone saw some fragments. Often the more accurate version was: knew a little, feared knowing more, therefore decided not to know.

\textbf{Resisters.} In February 1943, Munich University siblings Hans and Sophie Scholl and several companions formed the White Rose group, distributing antiwar, anti-Nazi leaflets. Arrested and executed four days later, they were only in their early twenties. Their leaflets broadly said silence was complicity and there was still time to tell the truth. In July 1944, Colonel Claus von Stauffenberg carried a bomb into Hitler's conference room, narrowly failing to assassinate him. Resisters existed. That fact matters: "there was no choice" was at least not a law of physics.

\textbf{A position that fits no box.} John Rabe was a German businessman and Nazi Party member living in Nanjing when Japanese forces captured it in 1937. Wearing a Nazi badge, he chaired the International Committee for the Nanking Safety Zone and used his German status to negotiate with Japanese forces. The zone sheltered more than two hundred thousand Chinese civilians. The label "Nazi Party member" did not predict his crucial choice. This does not absolve a party; it reminds us that position is not destiny. Decisions remain possible within it.

\textbf{Moscow: another version of the machine.} Nazis did not invent totalitarian politics alone. Under Stalin, the Soviet Great Purge of 1936--1938 brought public trials, coerced confessions, lists, and nighttime knocks. Hundreds of thousands were executed in two years; millions were sent to Gulag labor camps. Earlier, the Ukrainian famine of 1932--1933, produced by forced collectivization, killed millions. An ordinary Moscow engineer's situation strikingly resembled a Berlin clerk's: raise your hand in meetings, enthusiastically declare support, burn photographs after a neighbor's arrest, and consider every relative's name carefully when completing forms.

But differences among three machines matter. Mussolini's Italian Fascism, in power after the 1922 March on Rome, centered on \textbf{the supremacy of the state}: the state was everything, individuals nothing. It did not initially classify by race; systematic anti-Jewish laws arrived only in 1938 under Nazi pressure. Nazism added \textbf{biological racism} to Fascist state worship: ancestry determined everything, Jews were a "racial enemy" to be physically eliminated, and the "master race" needed living space. Stalinism came by another route, classifying through \textbf{class}, proclaiming proletarian liberation and targeting "class enemies" and "enemies of the people." The list could expand indefinitely: today's old Bolshevik became tomorrow's traitor. All three used one-party rule, secret police, leader worship, propaganda, and concentration camps, but justified hatred differently and imagined different ultimate ends. Grouping them as "totalitarianism" has been a scholarly tradition since the mid-twentieth century, represented by Arendt's The Origins of Totalitarianism. Other scholars have always resisted equating them. Different boundaries and victims genuinely affected who died and why.

Finally, give this opposing position its full reasons: "\textbf{Ordinary people had no choice.}" Listen before rebutting. First, fear was real: secret police threatened families as well as individuals; refusal could cost children's prospects or parents' pensions. Second, each step was small: today a sign, tomorrow a form, next day avoiding a neighborhood. No morning presented a choice labeled "do you support mass murder?" Third, information was monopolized: one radio voice, one newspaper vocabulary, an engineered information environment where seeing truth cost terribly much. Fourth, obedience was trained from childhood: schools, armies, and families had taught rule-following as virtue for decades. Fifth, livelihood was concrete: amid six million unemployed, a civil-service post meant the family ate. Every reason existed, and none is lightweight.

Set this beside the Scholls' story. Do not judge yet; remember the tension. Adequate reasons and "absolutely no choice" remain separated by a gap. The chapter will repeatedly measure its width.

\section{A Bridge for Thought: Five Gears}
The usual response to collective wrongdoing is moral astonishment: "How could they be so evil?" It comforts us by placing perpetrators outside "us," but explains nothing. If "evil" answers the question, Germany in 1933 somehow contained Europe's greatest concentration of evil people. That makes little sense.

Here is a portable tool: \textbf{replace "why were they so evil?" with "how does this machine run?"} Divide collective wrongdoing into five gears and inspect each.

\textbf{Gear one: propaganda.} Its chief function is not persuasion but \textbf{habituation}. Hear something a thousand times and you may not believe it, but it becomes "part of the environment," no longer worth rebutting. Its central craft is group portraiture: we are hardworking, injured, innocent; they are greedy, dangerous, not "us." A miniature example: by a rumor's tenth repost in a class chat, nobody remembers who first said it.

\textbf{Gear two: organization.} Its special skill is \textbf{fragmentation}. Divide a great undertaking into harmless-looking tasks: register, stamp, schedule, guard, collect keys. Nobody feels they performed the whole; each sees only a segment. Eichmann's defense matters because he articulated this mechanism himself. Its school equivalent: only one or two people strike during a group assault, but lookouts, messengers, hecklers, and video recorders each hold one small part of the chain.

\textbf{Gear three: fear.} It need not imprison everyone, only \textbf{make everyone see somebody imprisoned}. Punish a few and the majority censor themselves. Policing your own mouth costs less than any officer. The device travels across cultures. Wartime Japan organized ten-household neighborhood groups, distributing rations and monitoring thoughts through them. No Gestapo was necessary: neighborhood relations themselves became surveillance equipment.

\textbf{Gear four: interests.} Participation earns rewards; nonparticipation costs. Somebody moves into confiscated houses, fills vacant jobs, and takes the client list a "disappeared" Jewish colleague left on the desk. This gear's insidious feature is compatibility with conscience: someone can sincerely believe they merely "took over what nobody wanted."

\textbf{Gear five: obedience.} People are trained from childhood to obey authority. Handing responsibility to "orders" brings genuine relief: I no longer decide; I am just one link. We will devote a full section to this gear because it goes deeper than we like to admit.

The tool is simple. Whenever a news story, classroom, or group chat shows people jointly harming someone, suppress "how evil!" and ask: who repeats the slogans? Into what small tasks is the action divided? Who is publicly punished as an example? Who benefits? Who holds authority, and who receives obedience? Not every gear must be present, but once two or more mesh, "a few bad individuals" no longer explains the event.

This wrench serves more than history. It measures every machine that draws people in.

\section[Paper Acts Before the Street]{Reading Evidence: Paper Acts Before the Street}
Read two real documents, sixteen years apart. The first bills a defeated country; the second begins a killing process.

\textbf{First: Article 231 of the 1919 Treaty of Versailles.} The "war guilt clause" says, in a rendering of the Chinese paraphrase:

\begin{quote}
The Allied and Associated Governments affirm, and Germany acknowledges, that Germany and its allies bear responsibility for all losses and damage resulting from the war their aggression imposed on the Allied and Associated countries.
\end{quote}
(Source: Treaty of Versailles, Part VIII, Article 231, signed June 1919.)

Read its circumstances closely. This was not a neutral investigation but a "confession" the defeated had to sign: otherwise war continued. It placed not a conclusion but a thorn in German hands. For twenty years, that thorn was repeatedly exploited. Right-wing propaganda claimed defeat came not from losing militarily but from Jews and leftists "stabbing the country in the back." Reparations became enslavement; signatory Weimar politicians became "November criminals." These were not factual judgments but techniques for rubbing humiliation into fuel for hatred. A paper document does not kill by itself, but can prepare emotional foundations for killers.

\textbf{Second: Nazi Germany's Nuremberg Laws, September 1935.} The Reich Citizenship Law broadly restricted citizenship to people of "German or related blood," reducing Jews to state "subjects" without citizenship rights. The Law for the Protection of German Blood and German Honor, passed the same day, prohibited marriage between Jews and people of "German blood" and regulated even employment relationships in detail. (Source: the two laws enacted by the Reichstag at Nuremberg on September 15, 1935.)

Examine this stage. The killing assembly line began not with gas chambers but with \textbf{definition}. First rewrite a neighbor from "citizen" into "subject," engrave "us" and "them" in legal provisions, and confiscation, segregation, deportation, and extermination acquire "procedures." Clerks, accountants, and dispatchers now have forms. The machine's first screw turns on paper.

This paperwork had a thousand-year prehistory. Medieval European anti-Jewish persecution was religious: accusations of blasphemy, fabricated "blood ritual" rumors, expulsion from England in 1290 and Spain in 1492, and Venice's 1516 ghetto confining Jews to a designated neighborhood. In the nineteenth century, religious reasoning gave way to pseudoscience. The term Antisemitismus itself was coined in 1879 by Germany's Wilhelm Marr, treating Jewishness as "blood," not belief, so conversion changed nothing. In the early twentieth century, tsarist secret police forged The Protocols of the Elders of Zion, falsely alleging a Jewish plot for world rule; it circulated across Europe. In 1894, French officer Alfred Dreyfus was falsely accused of treason because he was Jewish, dividing France for over a decade. Nazis did not invent antisemitism. They connected accumulated religious hatred and nineteenth-century racial pseudoscience to a modern state: registration, railways, bureaucracy, propaganda industry. Old hatred was firewood; the modern state was the furnace.

Add a quotation that is not an exact quotation. German pastor Martin Niemöller initially supported Nazism, later entering a concentration camp for opposing state control of churches. After the war he repeatedly offered varying versions of this thought: first they arrested trade unionists, and I said nothing because I was not one; then Jews, and I said nothing because I was not Jewish; then Catholics, and I still said nothing; finally they came for me, and nobody remained to speak. It describes no atrocity directly, only the sequence of silences.

\section{One Machine, Three Explanations}
The facts are clear: a cultured country that produced Goethe, Kant, and Beethoven built an industrialized mass-murder machine in twelve years. Why? At least three serious explanations compete. Compare their explanatory reach, not their pleasing sound.

\textbf{First: Germany's special path.} Germany followed a different road from Britain and France: late unification, shallow democratic roots, powerful military and Junker elites, entrenched antisemitism and authoritarian obedience. Disaster grew from "Germanness." Its reasons are real: Weimar democracy was hastily built on defeat's ruins and pulled apart by far left and right. Its limits are equally evident. Antisemitism was not exclusively German: tsarist Russia produced the Protocols, France the Dreyfus affair. Weimar Berlin was also among Europe's freest, most creative cultural centers. The next counterexample is more damaging still.

\textbf{Counterexample one: a New Haven laboratory.} In 1961, American psychologist Stanley Milgram recruited ordinary Americans---workers, clerks, salesmen---near Yale in New Haven for an experiment supposedly studying punishment and learning. Participants played "teachers," pressing a shock button when a learner answered incorrectly. Voltage labels rose from 15 to 450 volts, with final settings marked "Danger: Severe Shock." The learner was actually in another room; screams were recorded. The real subject was the button-presser. Roughly two-thirds continued to maximum voltage. Some trembled, sweated, or asked to stop, but the white-coated experimenter calmly repeated "Please continue," "The experiment requires you to continue," and "We accept responsibility." Most kept pressing. Remember the location: not Munich, but Connecticut. If obedience were "German character," New Haven should not have produced this number.

\textbf{Second: the machine of modernity.} In Modernity and the Holocaust, sociologist Zygmunt Bauman offered a colder explanation: the Holocaust was not barbarism's brief counterattack against civilization but civilization's product. Without scientific racial classification, no lists of "who is Jewish"; without registers and archives, no locating them; without rail schedules and bureaucratic division of labor, no killing on this scale. Bureaucracy specializes in rewriting moral questions as technical ones: not "should it be done?" but "is it finished?" This powerfully explains the unprecedented scale. Yet it too faces a glaring counterexample.

\textbf{Counterexample two: Danish fishing boats.} Denmark also had modern bureaucracy, population registers, and ports. In October 1943, learning of an imminent Nazi roundup, people across society---police, doctors, fishers, taxi drivers---secretly carried over seven thousand Jews by small boat to neutral Sweden within weeks. More than ninety percent of Danish Jews survived the war. Bulgaria also resisted deportation pressure and protected its own Jews. Similar machinery produced opposite outcomes. Machinery may be necessary, but it does not turn itself. People operate it, and their society's answer to "who are we?" determines the gears' direction.

\textbf{Third: situations and human nature.} Milgram argued that ordinary people placed where "authority is present, responsibility handed upward, and demands gradually escalate" do things they would never do alone. The situation, not personality, commits the wrong. Decades of social psychology repeatedly tested and revised this explanation of "why good people cooperate." Its blind spot: reluctant laboratory button-pushers do not explain eager informers, fanatical volunteers, or people far exceeding orders. It also hides an ethical thorn. Attribute everything to situations and responsibility evaporates: everyone is situated, so nobody is responsible.

Each explanation holds a piece. Germany's historical conditions explain why the machine arose \textbf{there first}; modern bureaucracy explains why it \textbf{could} reach that scale; situational pressure explains why the gears \textbf{would} turn. Now pin this chapter's red line to the wall: \textbf{explaining mechanisms of participation never reduces participants' responsibility}. Understanding how machinery draws people in helps us recognize and dismantle it, not issue exemptions to those already inside. Nuremberg established a simple principle: obeying orders cannot excuse wrongdoing, because both those issuing and those executing orders are human.

\section[Is the Banality of Evil Enough?]{A Deeper Challenge: Is the "Banality of Evil" Enough?}
Now consider the chapter's weightiest theory and its controversy.

In the audience at Eichmann's 1961 Jerusalem trial sat Hannah Arendt, a Jewish political philosopher who had escaped Germany. Her magazine reports later became Eichmann in Jerusalem. Expecting a demon, she encountered something more disturbing: an ordinary-looking, ordinary-mannered man speaking administrative clichés, using "transport," "resettlement," and "Final Solution" for trainloads of lives. He stressed duty and lawfulness, even invoking Kant to justify obedience. Arendt proposed the famous concept of the \textbf{banality of evil}.

Define it carefully, because it is easily misunderstood. It does not mean evil is trivial or participants innocent. It means \textbf{enormous evil need not come from profound hatred; it may arise from "thoughtlessness"}. Not low intelligence, but refusing another's standpoint, replacing judgment with clichés and conscience with duty. Eichmann had a brain; he actively switched off its capacity to see from elsewhere. The theory removes a comforting sedative. If only demons do evil, we sleep peacefully, certain we are not demons. The "banality of evil" says wake up: you too stock the ingredients---obedience, clichés, unthinking behavior.

But honesty requires the story's second half. Later archival research complicated this portrait. Eichmann left extensive recorded conversations and manuscripts from exile in Argentina, recently assembled by scholars. There he was no "unthinking clerk" but a fanatical antisemite proud of slaughter. Performing ordinariness in court was a deliberate defense strategy. The "banality of evil" may therefore have misdiagnosed \textbf{Eichmann himself}.

Should the concept be discarded? No. Scholars broadly agree that it failed as a portrait of one man but remains sharp for \textbf{most positions within the machine}. Fanatical believers drew plans at Wannsee. Millions converting plans into schedules, forms, signs, and silence were largely Arendt's kind of people: not fanatical, merely unthinking, obedient, handing responsibility upward. A mass-murder machine needs two fuels: designers' hatred and cogs' banality. Without the first, it cannot be built; without the second, it cannot run.

Milgram supplied a measuring instrument. Later decades of repeated experiments revealed something more precise than "people obey." Participants generally continued not because they enjoyed harm---their distress is visible---but because they \textbf{could not find a face-saving way to refuse a calm, firm, supposedly legitimate authority who claimed responsibility}. Pressing a harmful button was easier than saying no to the white coat. This removes our last defense: we imagine we would stand up in that room, but two-thirds did not.

Combine theory and measurement and the chapter's central picture emerges. Totalitarian politics' greatest strength is not producing fanatics but \textbf{dividing the moral threshold for wrongdoing into five low steps}. Each is easy enough to cross without courage: grow used to a voice, post a sign, complete a form, lower your head, lower it again. At the final step, the owner of those footprints can no longer recognize the person who began.

Studying mechanisms has one purpose: recognize the staircase while its steps remain low. This does not reduce any Eichmann's sentence. On the contrary, precisely because each step is low, refusing each one can be demanded. That is what Nuremberg required.

\section{Today: The People's Receiver in Your Pocket}
Goebbels's loudmouth is gone. Your pocket contains a far more powerful People's Receiver.

Recommendation algorithms work on principles familiar to the old propaganda minister: repeat claims, profile groups, feed you what makes you pause. Anger is particularly good at stopping a scrolling finger. What the 1933 government needed years and millions of subsidized radios to build, today's platform builds the moment you register. This does not mean someone deliberately brainwashes you. It means designing information environments is no longer a state privilege but a standard phone feature. Identifying who repeats what and why you see it is basic self-protection today.

Run an online pile-on through the five gears. An unsourced screenshot is \textbf{propaganda}; by its tenth repost nobody remembers the origin. Participants divide tasks: exposing personal information, making memes, tagging schools and employers. Everyone holds one chain segment: \textbf{organization}. Defend the target and your screenshot comes next: \textbf{fear}. Followers, traffic, righteous satisfaction: \textbf{interests}. "Everybody says so" is \textbf{obedience} on a keyboard. What storm troopers once did door to door can happen overnight, each participant going to bed thinking they did almost nothing.

The school version is closer. In bullying, bullies and victims are minorities; most people stand around them, laughing along, pretending not to see, later saying they too were uncomfortable. Bystanding is not neutrality but oxygen for this small machine. A wrench jamming the gears may be surprisingly small: "That's enough," sitting beside someone isolated, refusing a repost. History does not require everyone to be a Scholl. It proves many positions exist between zero and one hundred.

Zoom out to other clothes and continents. In Rwanda in 1994, a radio station repeatedly called Tutsi people "cockroaches." From April to July, about eight hundred thousand were killed in roughly one hundred days, many by machetes, neighbors killing neighbors. Radio repeated its German work among equatorial hills. Earlier, in 1915, Armenians in the Ottoman Empire were deported into the Syrian desert in groups, about a million dying en route, through the familiar bureaucratic sequence of registration, confiscation, deportation. Nor do these gears require totalitarian government. In 1950s America, without secret police or one-party rule, McCarthyism used public "are you or aren't you?" questioning, loyalty checks, and blacklists to cost thousands of teachers, actors, and civil servants jobs and voices. Democracy is an immune system, not a vaccine, and immune systems can fail.

This is not an invitation to cry "Nazis!" at everything. Precisely because real totalitarian machinery is so weighty, the word should not be squandered. Tools are for measurement. Each gear has a scale: measure before speaking.

\section{Your Question: Where Is the Line?}
Return to the radio and the lowered heads before the bakery window.

Answer the chapter's final question, with reasons: \textbf{up to what point is obedience understandable, and from what point unforgivable?} If you drew the line for residents of that street in 1935, where would it fall?

Weigh your evidence once more before drawing.

On one side: fear and threats to family were real. Each step was small; no morning announced "wrongdoing begins today." Amid six million unemployed, a job fed a family. Information monopoly made truth terribly costly. People were trained to obey, and two-thirds in Milgram's room went all the way. You cannot be certain you belong to the other third.

On the other: Danish boats were real, saving over seven thousand lives within weeks. White Rose leaflets were real, costing two young adults their heads. Rabe in Nanjing was real; a label did not decide for him. Warsaw's buried metal containers were real; people kept recording in the machine's depths. Each proves "no choice" is not a physical law. The gap between "adequate reasons" and "no choice" exists.

The heaviest weight is our red line: understand recruitment into machinery to dismantle it, not to excuse its components. Yet if refusing the next step can be demanded, where does that demand begin? Posting the sign? Filling the form? Lowering your head in line? Earlier, when the radio first called neighbors vermin and you did not change stations?

Where will you draw your line, and what reasons will defend it? If someone posts the first sign tomorrow in your chat, classroom, or building, will those reasons be enough to make you raise your head?
